% ==============================================================================
% Paquetes y comandos definidos por el usuario
% ==============================================================================
\usepackage{multirow}
\usepackage{makecell}
\usepackage{tabularx}
\usepackage{colortbl}
\usepackage{listings}
\usepackage[most]{tcolorbox}
\usetikzlibrary{babel}
\usepackage{pmboxdraw}
\usepackage{fontawesome5}

% Las comillas dentro del código se escriben tal cual (babel en español las usa
% como atajo para diéresis)
\AtBeginDocument{\shorthandoff{"}}

% Caracteres Unicode que pdflatex no trae por defecto
\DeclareUnicodeCharacter{1F427}{\faLinux}
\DeclareUnicodeCharacter{1F34E}{\faApple}
\DeclareUnicodeCharacter{2753}{\textcolor{red}{\faQuestion}}
\DeclareUnicodeCharacter{2600}{\faSun}
\DeclareUnicodeCharacter{FE0F}{}

% El código en línea no se separa en sílabas; esto evita que desborde el margen
\setlength{\emergencystretch}{3em}
% Permite cortar las URL de la bibliografía en cualquier carácter
\setcounter{biburllcpenalty}{7000}
\setcounter{biburlucpenalty}{8000}
\setcounter{biburlnumpenalty}{9000}

% Columna X de tabularx alineada a la izquierda
\newcolumntype{L}{>{\raggedright\arraybackslash}X}

% ------------------------------------------------------------------------------
% Bloques de código
% ------------------------------------------------------------------------------
\definecolor{fondo-codigo}{gray}{0.97}
\definecolor{fondo-caja}{gray}{0.90}
\definecolor{codigo-clave}{HTML}{D73A49}
\definecolor{codigo-cadena}{HTML}{298E0D}
\definecolor{codigo-comentario}{HTML}{8A8A8A}

\lstdefinelanguage{nix}{
    morekeywords={let,in,if,then,else,with,inherit,rec,import,assert,or,true,false,null},
    sensitive=true,
    morecomment=[l]{\#},
    morecomment=[s][\color{codigo-cadena}]{''}{''},
    morestring=[b]",
}

\lstdefinelanguage{TypeScript}{
    morekeywords={import,from,export,default,const,let,var,new,return,function,
        if,else,for,while,true,false,null,undefined,type,interface,class,extends,
        typeof,as,string,number,boolean},
    sensitive=true,
    morecomment=[l]{//},
    morecomment=[s]{/*}{*/},
    morecomment=[s][\color{codigo-cadena}]{`}{`},
    morestring=[b]",
    morestring=[b]',
}

\lstset{
    basicstyle=\ttfamily\footnotesize,
    keywordstyle=\color{codigo-clave},
    stringstyle=\color{codigo-cadena},
    commentstyle=\color{codigo-comentario},
    columns=fixed,
    basewidth=0.525em,
    keepspaces=true,
    showstringspaces=false,
    breaklines=true,
    upquote=true,
    aboveskip=0pt,
    belowskip=0pt,
    literate=
        {á}{{\'a}}1 {é}{{\'e}}1 {í}{{\'i}}1 {ó}{{\'o}}1 {ú}{{\'u}}1
        {Á}{{\'A}}1 {É}{{\'E}}1 {Í}{{\'I}}1 {Ó}{{\'O}}1 {Ú}{{\'U}}1
        {ñ}{{\~n}}1 {Ñ}{{\~N}}1 {ü}{{\"u}}1
        {¿}{{\textquestiondown}}1 {¡}{{\textexclamdown}}1
        {°}{{\textdegree}}1 {—}{{\textemdash}}1 {–}{{\textendash}}1
        {─}{{\textSFx}}1 {│}{{\textSFxi}}1 {└}{{\textSFii}}1 {├}{{\textSFviii}}1
        {☀}{{\faSun}}1 {️}{}0,
}

% Bloque de código con fondo gris
\newtcolorbox{bloquecodigo}{
    enhanced, breakable, colback=fondo-codigo, frame hidden, arc=4pt,
    boxsep=0pt, left=12pt, right=12pt, top=12pt, bottom=12pt,
}

% Opción de selección múltiple: etiqueta a la izquierda del bloque de código
\newtcolorbox{opcion}[1]{
    enhanced, colback=fondo-codigo, frame hidden, arc=4pt,
    boxsep=0pt, left=12pt, right=12pt, top=12pt, bottom=12pt,
    enlarge left by=2em, width=\linewidth-2em,
    overlay={\node[anchor=north east, inner sep=0pt, font=\bfseries]
        at ([xshift=-1em]frame.north west) {#1};},
}

% Recuadro gris para textos que se leen al participante
\newtcolorbox{cajagris}{
    enhanced, colback=fondo-caja, frame hidden, arc=4pt,
    boxsep=0pt, left=12pt, right=12pt, top=12pt, bottom=12pt,
}

% ------------------------------------------------------------------------------
% Instrumentos de los anexos
% ------------------------------------------------------------------------------
% Casilla vacía para marcar
\newcommand{\cuadrito}{{\setlength{\fboxsep}{0pt}\fbox{\rule{0pt}{0.45cm}\rule{0.45cm}{0pt}}}}
% Casilla con su etiqueta
\newcommand{\casilla}[1]{\cuadrito\hspace{0.4cm}#1\quad}
% Línea para escribir a mano
\newcommand{\lineafirma}{\rule[-1pt]{0.4\linewidth}{0.5pt}}
